% !TEX root = ../main.tex
% Section: Appendix: Counting Covariance-Level Mean-Field Laws

\section{Counting Covariance-Level Mean-Field Laws}
\label{app:law-counting}

This appendix records the Lie-algebra computation behind the covariance-level
completeness statement used in Section~\ref{sec:laws}.
It is a continuum, covariance-level analogue of the local/generic finite-width
completeness theorem of \citet{MarcotteGribonvalPeyre2023}.

Width-consistent laws and their Wasserstein interpretation are defined
in Appendix~\ref{app:mean-field}. We work here directly with functions
of the covariance and the vector fields that must annihilate them.

\subsection{Split Bures vector fields}
\label{subsec:split-bures-fields}

To characterize the laws shared by the architecture, we allow every
smooth split loss:
\[
F(\Sigma):=\ell(\Svu),
\qquad
\Sigma=
\begin{pmatrix}
\Suu&\Suv\\
\Svu&\Svv
\end{pmatrix}.
\]
If \(E:=\nabla \ell(\Svu)\), set
\[
G_E:=
\begin{pmatrix}
0&E^\top\\
E&0
\end{pmatrix}\in\mathfrak p_J,
\qquad
\mathfrak p_J:=\{G\in\Ssym^m:GJ+JG=0\}.
\]
The corresponding Bures vector field is
\begin{equation}
\label{eq:appendix-split-vector-field}
X_\ell(\Sigma):=-(G_{\nabla \ell(\Svu)}\Sigma+\Sigma G_{\nabla \ell(\Svu)}).
\end{equation}
We denote the resulting space of split Bures vector fields by
\[
\mathcal X_{\rm split}
:=
\{X_\ell:\ell\in C^\infty(\R^{d_v\times d_u})\}.
\]
It is a real vector space.  Its constant-gradient subspace is obtained from
linear losses \(\ell_E(\Svu):=\ip{E}{\Svu}\):
\[
\mathcal X_{\rm lin}
:=
\{X_E:E\in\R^{d_v\times d_u}\},
\qquad
X_E(\Sigma):=-(G_E\Sigma+\Sigma G_E).
\]

\subsection{Generated Lie algebra}

Lie brackets reveal the directions generated by all split losses,
which determine how many independent common laws can remain.
Let
\[
\mathfrak o(J):=\{A\in\R^{m\times m}:A^\top J+JA=0\}
\]
be the Lie algebra of \(O(d_u,d_v)\).  For \(A\in\mathfrak o(J)\), define the
infinitesimal congruence field
\[
\xi_A(\Sigma):=-(A\Sigma+\Sigma A^\top).
\]
Its flow is \(\Sigma(t)=e^{-tA}\Sigma(0)e^{-tA^\top}\), so it is tangent to
the \(O(d_u,d_v)\)-congruence orbit of \(\Sigma\).  When \(A=G_E\in\mathfrak p_J\),
\(\xi_A=X_E\).

\begin{proposition}[Split fields generate \(\mathfrak o(d_u,d_v)\)]
\label{prop:appendix-lie-generated}
With the bracket convention \([X,Y]:=DY[X]-DX[Y]\), the map
\[
A\longmapsto \xi_A
\]
is a Lie-algebra homomorphism:
\[
[\xi_A,\xi_B]=\xi_{[A,B]}.
\]
Moreover,
\[
\operatorname{Lie}(\mathfrak p_J)=\mathfrak o(J).
\]
Consequently the Lie algebra generated by
\(\mathcal X_{\rm lin}\), the constant-gradient split Bures fields, is
\[
\{\xi_A:A\in\mathfrak o(J)\}.
\]
\end{proposition}

\begin{proof}
The homomorphism identity follows by differentiating the linear vector fields
\(\xi_A\).  It remains to compute the matrix Lie algebra generated by
\(\mathfrak p_J\).  For
\[
G_E=
\begin{pmatrix}
0&E^\top\\
E&0
\end{pmatrix},
\qquad
G_F=
\begin{pmatrix}
0&F^\top\\
F&0
\end{pmatrix},
\]
one has
\[
[G_E,G_F]
=
\begin{pmatrix}
E^\top F-F^\top E&0\\
0&EF^\top-FE^\top
\end{pmatrix}.
\]
These brackets span the block-diagonal skew-symmetric algebra
\[
\mathfrak k_J:=
\left\{
\begin{pmatrix}
A&0\\
0&D
\end{pmatrix}:A^\top=-A,\ D^\top=-D
\right\},
\]
by choosing \(E,F\) supported on one common row or one common column.  Since
\[
\mathfrak o(J)=\mathfrak k_J\oplus\mathfrak p_J,
\]
the generated Lie algebra is all of \(\mathfrak o(J)\).
\end{proof}

\paragraph{Fields versus their generated distribution.}
The full space \(\mathcal X_{\rm split}\) contains nonlinear fields and
is not identified with a finite-dimensional matrix Lie algebra. Each
of its fields is a smooth, state-dependent linear combination of a
basis of the \(X_E\)'s. Consequently all such fields and their brackets
remain tangent to the same congruence orbits, while the constant-gradient
fields already generate the whole orbit tangent space. It is this
distribution that determines the common conservation laws.

\subsection{Completeness of the spectral laws}

The generated directions coincide with the tangent directions of
covariance congruence orbits; their codimension gives the local law count.
For a constant-gradient split Bures field, \(B:=J\Sigma\) satisfies
\[
\dot B=[G_E,B].
\]
The trace powers $\mathcal H_k$ in \eqref{eq:bures-G} are
therefore conserved. The next corollary states that, locally and generically, these
are the only smooth covariance-level laws common to all split losses.

Let
\[
\Omega_{\rm simp}
:=
\{\Sigma\in\Spp^m:J\Sigma\text{ has simple spectrum}\}.
\]

\begin{corollary}[Marcotte-style covariance-level completeness]
\label{cor:appendix-covariance-completeness}
Let \(h\) be a smooth function on an open subset of \(\Omega_{\rm simp}\).
If \(h\) is conserved by every split Bures vector field \(X_\ell\) of
\eqref{eq:appendix-split-vector-field}, then \(h\) is locally a smooth
function of the eigenvalues of \(J\Sigma\), equivalently of
\[
\mathcal H_1(\Sigma),\ldots,\mathcal H_m(\Sigma).
\]
These \(m\) laws are functionally independent on \(\Omega_{\rm simp}\).
Thus the characteristic polynomial of \(J\Sigma\) gives a complete
local/generic list of covariance-level mean-field conservation laws for the
split two-layer linear architecture.
\end{corollary}

\begin{proof}
It is enough to use the linear split losses.  If \(X_Eh=0\) for every \(E\),
then \(h\) is also annihilated by all iterated Lie brackets of the \(X_E\)'s.
By Proposition~\ref{prop:appendix-lie-generated},
\[
\mathrm{d}h(\Sigma)[\xi_A(\Sigma)]=0
\qquad\text{for all }A\in\mathfrak o(J).
\]
Thus \(h\) is constant on the connected pieces of the
\(O(d_u,d_v)\)-congruence orbits.

The orbit invariants are exactly the eigenvalues of \(J\Sigma\) on
\(\Omega_{\rm simp}\).  Indeed,
\[
J(Q\Sigma Q^\top)=Q^{-\top}(J\Sigma)Q^\top
\qquad(Q^\top JQ=J),
\]
so the spectrum is preserved.  Conversely, for \(\Sigma\succ0\), the matrix
\(J\Sigma\) is similar to the symmetric matrix
\(\Sigma^{1/2}J\Sigma^{1/2}\), hence has real nonzero eigenvalues with
\(d_u\) positive and \(d_v\) negative signs.  On the simple-spectrum stratum
its \(J\)-orthogonal eigenbasis can be locally normalized to a
\(J\)-orthonormal basis, so equal locally ordered eigenvalues imply membership
in the same local \(O(d_u,d_v)\)-orbit.  Therefore any smooth orbit-constant
function factors locally through the eigenvalues.  Newton identities replace
the eigenvalues by \(\mathcal H_1,\ldots,\mathcal H_m\).
More explicitly, for the locally ordered eigenvalues \(\lambda_j\),
\[
 \frac{\partial\mathcal H_k}{\partial\lambda_j}
 =k\lambda_j^{k-1},\qquad 1\le k,j\le m.
\]
This scaled Vandermonde matrix is invertible at distinct eigenvalues.
The eigenvalues themselves vary independently by varying the diagonal
covariance in a fixed normalized eigenbasis. Hence the spectral map has
rank \(m\), and its local level sets have dimension
\(m(m+1)/2-m=m(m-1)/2\), exactly the orbit dimension.
\end{proof}

For empirical covariances \(\Sigma=WW^\top/n\) with
\(W=\binom{U}{V}\), the finite-width Gram imbalance is
\[
\Gamma:=V^\top V-U^\top U=-W^\top JW.
\]
The nonzero spectra of \(J\Sigma=JWW^\top/n\) and
\(-\Gamma/n=W^\top JW/n\) coincide.  Corollary~\ref{cor:appendix-covariance-completeness}
is therefore the covariance quotient of the Marcotte--Gribonval--Peyr\'e
completeness theorem: the full labeled finite-width conservation data include the matrix
\(\Gamma\), while the covariance-level law keeps only its spectrum.

\subsection{A ReLU contrast}

Homogeneity preserves a scalar balance beyond linear networks, but it
does not give the same covariance closure or spectral family.
For a two-layer ReLU or more generally one-homogeneous network, the characteristic
equations associated with a potential \(v^\top\psi(u)\), with \(\psi\) one-homogeneous
and possibly time-dependent,
have the form
\[
\dot u=-D\psi(u)^\top v,\qquad \dot v=-\psi(u).
\]
Euler's identity \(D\psi(u)u=\psi(u)\) gives
\[
\frac{\mathrm{d}}{\mathrm{d}t}\bigl(\norm u^2-\norm v^2\bigr)=0.
\]
For ReLU this calculation is interpreted almost everywhere along the
characteristics, with an activation-boundary convention satisfying Euler's
identity (for instance derivative zero at a zero preactivation).
Thus the full mean-field flow preserves the pushforward distribution of
\(\norm u^2-\norm v^2\), and its covariance-level shadow is the total imbalance
\[
\int\bigl(\norm u^2-\norm v^2\bigr)\,\mathrm{d}\mu(u,v).
\]
Unlike the linear split case above, this ReLU statement is not obtained from a
closed finite-dimensional Bures covariance equation; it is included only to
highlight how much smaller the covariance-level conservation-law problem can be
than the full particle-level one. No completeness claim for ReLU laws is
made here.

%%% Local Variables:
%%% mode: LaTeX
%%% TeX-master: "../main"
%%% End:
